\documentclass[11pt]{article}
\usepackage[margin=1in]{geometry}
\usepackage{amsmath, amssymb}
\usepackage{booktabs}
\usepackage{multirow}
\usepackage{xcolor}
\usepackage{hyperref}
\usepackage{microtype}

\title{Loss Function Comparison for Privileged Self-Distillation}
\author{OPSD Project}
\date{\today}

\begin{document}
\maketitle

\section{Setup}

All runs use \textbf{Qwen3-1.7B} with LoRA ($r=64$, $\alpha=128$, all projection layers), trained for 200 steps on OpenThoughts-Math-30k. The teacher has access to privileged information (ground-truth solution in context); the student sees only the problem. The teacher branch is frozen via the fixed-teacher strategy: LoRA adapters update only the student forward pass, while the teacher uses the frozen base model weights.

\noindent\textbf{Shared hyperparameters:}
\begin{itemize}
  \item Learning rate: $5\times10^{-6}$, max grad norm $0.1$
  \item Batch: 4 per device $\times$ 4 GPUs $\times$ 2 gradient accumulation steps $= 32$ effective
  \item Max completion length: 1024 tokens; temperature $\tau = 1.1$, top-$p = 0.95$, top-$k = 20$
  \item On-policy generation via vLLM (colocate mode, 4$\times$L40S GPUs)
\end{itemize}

\section{Loss Functions}

Let $q_\theta(y_t \mid y_{<t}, x)$ be the student and $p(y_t \mid y_{<t}, x, c)$ the frozen privileged teacher. All losses operate on sampled tokens only (no full vocabulary summation).

\subsection{Token-KL / Reverse KL (revkl)}

Per-token reverse KL used as a policy-gradient surrogate~\cite{thinkingmachines2025}:
\begin{equation}
  r_t = \log p(y_t \mid y_{<t}, x, c) - \log q_\theta(y_t \mid y_{<t}, x)
\end{equation}
\begin{equation}
  \mathcal{L}_{\text{revkl}} = -\frac{1}{T}\sum_{t=1}^T \operatorname{sg}(r_t) \cdot \log q_\theta(y_t \mid y_{<t}, x)
\end{equation}
Credit assignment is purely local: each token is weighted by its own log-ratio with no information from the rest of the sequence.

\subsection{Cumulative Sequence KL (cumseq)}

Inspired by \citet{cumseq2025}, this adds long-range credit assignment via the suffix mean log-ratio:
\begin{equation}
  \tilde{R}_t = \frac{1}{T - t + 1}\sum_{s=t}^{T} r_s
\end{equation}
\begin{equation}
  \mathcal{L}_{\text{cumseq}} = -\frac{1}{T}\sum_{t=1}^T \operatorname{sg}(\tilde{R}_t) \cdot \log q_\theta(y_t \mid y_{<t}, x)
\end{equation}
The suffix sum is divided by remaining valid token count at each position, bounding the weight to $[\min_t r_t,\, \max_t r_t]$ and preventing the $\mathcal{O}(T)$ magnitude blowup that arises from unnormalized suffix sums at long completion lengths.

\subsection{MiniLLM-Style Loss (minillm)}

The exact policy-gradient form derived by \citet{minillm2023} for minimizing $\mathrm{KL}(q_\theta \| p)$ at the sequence level:
\begin{equation}
  \nabla_\theta \mathcal{L}_{\mathrm{KL}} = -\mathbb{E}_{y \sim q_\theta}\left[\sum_{t=1}^T R_{t-1} \cdot \nabla_\theta \log q_\theta(y_t \mid y_{<t}, x)\right]
\end{equation}
where $R_{t-1} = \sum_{s \geq t} r_s$ is the return-to-go from position $t$ (inclusive). We apply the same suffix-mean normalization as cumseq for stability, then subtract the $-1$ baseline that appears in the MiniLLM gradient derivation:
\begin{equation}
  w_t = \operatorname{sg}(\tilde{R}_t) - 1
\end{equation}
\begin{equation}
  \mathcal{L}_{\text{minillm}} = -\frac{1}{\sum_t \mathbf{1}_t}\sum_{t=1}^T w_t \cdot \log q_\theta(y_t \mid y_{<t}, x)
\end{equation}
The $-1$ term shifts the weight so that a flat suffix ($\tilde{R}_t = 1$) produces zero gradient, reducing variance relative to cumseq. The normalization is also global over tokens rather than per-sequence.

\paragraph{Relationship between losses.}
The gradient of $\mathcal{L}_{\text{cumseq}}$ is $\tilde{R}_t \cdot \nabla_\theta \log q_\theta(y_t)$, while that of $\mathcal{L}_{\text{minillm}}$ is $(\tilde{R}_t - 1) \cdot \nabla_\theta \log q_\theta(y_t)$. The difference is a constant additive term $\nabla_\theta \log q_\theta(y_t)$ — an entropy-like bonus present in cumseq but absent in minillm. The revkl gradient is $r_t \cdot \nabla_\theta \log q_\theta(y_t)$, i.e.\ cumseq without long-range credit assignment.

\section{Evaluation}

Evaluation uses vLLM with thinking mode enabled (\texttt{/think} appended to prompts, Qwen3 style). AIME problems use 32k max generation tokens; all other datasets use 16k. Accuracy is measured by \texttt{math\_verify} exact-match on $\backslash$\texttt{boxed\{\}} answers.

\subsection{Greedy Accuracy (pass@1, thinking)}

\begin{table}[h]
\centering
\caption{Greedy accuracy (\%) with thinking mode. AIME: 30 problems each. Others: 300 samples each. Pending entries marked with \textemdash.}
\label{tab:greedy}
\begin{tabular}{lcccccccc}
\toprule
\multirow{2}{*}{\textbf{Loss}} & \multicolumn{2}{c}{\textbf{AIME}} & \multirow{2}{*}{\textbf{AMC}} & \multirow{2}{*}{\textbf{GSM8K}} & \multirow{2}{*}{\textbf{MATH}} & \multirow{2}{*}{\textbf{Minerva}} & \multirow{2}{*}{\textbf{Olympiad}} & \multirow{2}{*}{\textbf{Other Avg}} \\
\cmidrule(lr){2-3}
 & \textbf{24} & \textbf{25} & & & & & & \\
\midrule
Base (Qwen3-1.7B) & 46.7 & 23.3 & 67.5 & 90.3 & 72.7 & 23.5 & 48.7 & 60.2 \\
\midrule
revkl             & 40.0 & 26.7 & \textbf{71.1} & 90.0 & 72.0 & 22.8 & \textbf{49.0} & 60.1 \\
cumseq            & --- & --- & 63.9 & 89.0 & 72.0 & \textbf{23.9} & 48.3 & 59.4 \\
minillm           & --- & --- & --- & --- & --- & --- & --- & --- \\
\bottomrule
\end{tabular}
\end{table}

\subsection{Pass@12 (sampling, thinking)}

Sampling uses temperature $0.6$, top-$p = 0.95$, top-$k = 20$, $k=12$ candidates per problem.

\begin{table}[h]
\centering
\caption{Pass@12 (\%) with thinking mode. All entries pending.}
\label{tab:pass12}
\begin{tabular}{lcccccccc}
\toprule
\multirow{2}{*}{\textbf{Loss}} & \multicolumn{2}{c}{\textbf{AIME}} & \multirow{2}{*}{\textbf{AMC}} & \multirow{2}{*}{\textbf{GSM8K}} & \multirow{2}{*}{\textbf{MATH}} & \multirow{2}{*}{\textbf{Minerva}} & \multirow{2}{*}{\textbf{Olympiad}} & \multirow{2}{*}{\textbf{Other Avg}} \\
\cmidrule(lr){2-3}
 & \textbf{24} & \textbf{25} & & & & & & \\
\midrule
Base (Qwen3-1.7B) & --- & --- & --- & --- & --- & --- & --- & --- \\
revkl             & --- & --- & --- & --- & --- & --- & --- & --- \\
cumseq            & --- & --- & --- & --- & --- & --- & --- & --- \\
minillm           & --- & --- & --- & --- & --- & --- & --- & --- \\
\bottomrule
\end{tabular}
\end{table}

\section{Notes}

\begin{itemize}
  \item The cumseq run used a suffix-\emph{sum} (unnormalized) in its original implementation, causing $\mathcal{O}(T)$ loss magnitude growth and training divergence. The suffix-\emph{mean} normalization was applied as a fix prior to the run whose results appear above.
  \item The minillm run completed training (200 steps, final loss $-0.03$, grad norm $0.42$); evaluation jobs are pending.
  \item AIME results show high variance due to small test set size ($n=30$); differences of $\pm 1$ problem correspond to $\pm 3.3\%$.
  \item numina\_eval omitted from Other Avg above as it was not included in the current evaluation sweep.
\end{itemize}

\bibliographystyle{plain}
\begin{thebibliography}{9}

\bibitem{minillm2023}
Gu et al.\ (2023).
\textit{MiniLLM: Knowledge Distillation of Large Language Models}.
arXiv:2306.08543.

\bibitem{cumseq2025}
Anonymous (2025).
\textit{Cumulative Sequence KL Distillation}.
arXiv:2506.09477.

\bibitem{thinkingmachines2025}
Thinking Machines Lab (2025).
\textit{On-Policy Reverse KL Distillation}.
\url{https://thinkingmachines.ai}.

\end{thebibliography}

\end{document}
